\author{Brett Reynolds}
\title{언어의 풍경: Language Landscapes}
\subtitle{영어 교사를 위한 영어 구조 안내서[한국어판]}
% \ISBNdigital{}
% \ISBNhardcover{}
% \BookDOI{}
\typesetter{Brett Reynolds}
% \proofreader{}
% \lsCoverTitleSizes{51.5pt}{17pt}
\renewcommand{\lsSeries}{tbls}
%\newcommand{\lsSeriesTitle}{Textbooks in Language Sciences}
%\newcommand{\lsSeriesColor}{lsYellow}
\renewcommand{\lsSeriesNumber}{16}
\BackBody{언어는 외워야 할 목록이 아니라 탐구해야 할 풍경입니다.

좋은 현장 자연학자는 이름만 붙이지 않습니다. 무엇이 어떻게 행동하는지, 어디에 모이는지, 왜 달라지는지를 관찰합니다. 이 책은 영어를 같은 방식으로 읽도록 안내합니다. 규칙서가 영어가 어떠해야 한다고 말하는지가 아니라, 실제 화자가 무엇을 하는지에 주목하는 것입니다. 이론적 틀은 \textit{The Cambridge Grammar of the English Language}에 바탕을 두되, 접근하기 쉽게 단순화했습니다. 하지만 수준을 낮추지는 않았습니다.

각 장은 소리와 단어에서 구, 절, 담화로 나아가며, 분명한 영어 예문을 다른 언어의 단면과 함께 제시합니다. 각 장에서 같은 과정을 연습합니다. 패턴을 관찰하고, 증거로 검토하고, 어디에서 그 패턴이 무너지는지 보고, 그 무너짐이 무엇을 알려 주는지 묻는 과정입니다. 끝까지 읽고 나면, 교재를 평가하고, 학습자의 문장이 왜 조금 어색한지 설명하고, 질문을 받을 때 자신의 설명을 근거 있게 방어할 수 있는 실용적인 문법 지식을 갖게 될 것입니다.

교사 양성 과정에 있든 이미 교실에서 가르치고 있든, 얻는 것은 같습니다. 정당화할 수 없는 규칙에 기대는 일을 멈추고, 언어 체계를 스스로 볼 수 있게 됩니다.}
